\section*{Capítulo \ref{ch:b1:crystals-ionic-covalent} --- Cristais II: sólidos iônicos, covalentes e moleculares}

\begin{solution}{exo:b1:crystals-ionic-covalent:1}
Ânions: $8 \times \frac18 + 6 \times \frac12 = 4$; cátions: $12 \times
\frac14 + 1 = 4$. Cada íon tem 6 vizinhos do outro tipo (6:6). Quatro
cátions e quatro ânions: \ce{NaCl}.
\end{solution}

\begin{solution}{exo:b1:crystals-ionic-covalent:2}
Metade da diagonal principal: $a\sqrt3/2 = 412.3 \times 0.866 = \qty{357.1}{pm}$.
Um \ce{Cs+} e $8 \times \frac18 = 1$ \ce{Cl-}: uma unidade de fórmula.
\end{solution}

\begin{solution}{exo:b1:crystals-ionic-covalent:3}
$x = 102/181 = 0.56$, entre 0.414 e 0.732: coordenação octaédrica
(6), o \omterm{def:b1:crystals-ionic-covalent:structure-type}{tipo sal-gema}.
\end{solution}

\begin{solution}{exo:b1:crystals-ionic-covalent:4}
Metálico: cobre. Iônicos: cloreto de sódio, fluoreto de cálcio.
Covalentes: diamante, quartzo. Moleculares: di-iodo, gelo.
\end{solution}

\begin{solution}{exo:b1:crystals-ionic-covalent:5}
Ao longo de uma aresta, ânion--cátion--ânion em contato: $a = 2(r_+ + r_-)$.
Ao longo de uma diagonal de face ficam dois ânions à distância
$a\sqrt2/2$; eles não podem se interpenetrar: $a\sqrt2/2 \ge 2r_-$, logo
$(r_+ + r_-)\sqrt2 \ge 2r_-$, ou seja,
$r_+/r_- \ge \sqrt2 - 1$.
\end{solution}

\begin{solution}{exo:b1:crystals-ionic-covalent:6}
$\rho = 4 \times 58.44/(6.022 \times 10^{23} \times (5.641 \times
10^{-8})^3) = \qty{2.16}{g/cm^3}$, contra \qty{2.17}{g/cm^3} medidos.
\end{solution}

\begin{solution}{exo:b1:crystals-ionic-covalent:7}
$\rho = 168.36/(6.022 \times 10^{23} \times (4.123 \times 10^{-8})^3) =
\qty{3.99}{g/cm^3}$, o valor medido.
\end{solution}

\begin{solution}{exo:b1:crystals-ionic-covalent:8}
\ce{Zn-S} $= a\sqrt3/4 = \qty{234.2}{pm}$;
$\rho = 4 \times 97.44/(6.022 \times 10^{23} \times (5.409 \times
10^{-8})^3) = \qty{4.09}{g/cm^3}$, próximo do valor medido,
\qty{4.04}{g/cm^3}, da esfalerita natural.
\end{solution}

\begin{solution}{exo:b1:crystals-ionic-covalent:9}
\ce{C-C} $= 245.6/\sqrt3 = \qty{141.8}{pm}$; entre camadas, $669.6/2 =
\qty{334.8}{pm}$. Volume da célula: $\frac{\sqrt3}{2}a^2c = \qty{3.498e-23}{cm^3}$,
$\rho = 4 \times 12.01/(6.022 \times 10^{23} \times 3.498 \times 10^{-23})
= \qty{2.28}{g/cm^3}$, contra \qty{3.52}{g/cm^3} para o diamante: as
camadas estão longe umas das outras, mantidas apenas por forças de London.
\end{solution}

\begin{solution}{exo:b1:crystals-ionic-covalent:10}
Veja a demonstração da \cref{prop:b1:crystals-ionic-covalent:diamond}:
$C = \pi\sqrt3/16 = 0.34$. A dureza vem da força e da rigidez
tridimensional das \omterm{def:b1:lewis-resonance-vsepr:lewis}{ligações covalentes}, não do empacotamento: cada átomo
tem apenas quatro vizinhos, em ângulos fixos, o que deixa muito espaço
vazio.
\end{solution}

\begin{solution}{exo:b1:crystals-ionic-covalent:11}
\ce{Si-Si} $= a\sqrt3/4 = \qty{235.2}{pm}$, \omterm{def:b1:periodicity:radius}{raio covalente}
\qty{117.6}{pm}; $\rho = 8 \times 28.09/(6.022 \times 10^{23} \times
(5.431 \times 10^{-8})^3) = \qty{2.33}{g/cm^3}$, o valor medido.
\end{solution}

\begin{solution}{exo:b1:crystals-ionic-covalent:12}
$x = 152/181 = 0.84 > 0.732$: a regra prevê o \omterm{def:b1:crystals-ionic-covalent:structure-type}{tipo cloreto de césio}.
Na estrutura do sal-gema, \ce{Rb-Cl} $= a/2 = \qty{329.1}{pm}$, próximo de
$152 + 181 = \qty{333}{pm}$: a estrutura observada é coerente. A regra
trata os íons como esferas rígidas e ignora sua \omterm{def:b1:periodicity:polarisability}{polarizabilidade} e a
pequena diferença de energia entre as duas estruturas (o cloreto de
rubídio adota o \omterm{def:b1:crystals-ionic-covalent:structure-type}{tipo cloreto de césio} sob pressão); é um guia, não uma
lei.
\end{solution}

\begin{solution}{pb:b1:crystals-ionic-covalent:1}
\textbf{1.} Nos 8 vértices e nos 6 centros de face: $8 \times \frac18 + 6
\times \frac12 = 4$ íons.
\textbf{2.} 8 \omterm{def:b1:crystals-metals:site}{sítios tetraédricos}, todos ocupados: 8 \ce{F-} para 4
\ce{Ca^{2+}}, razão 1:2, \ce{CaF2}.
\textbf{3.} \ce{F-}: 4 \ce{Ca^{2+}} nos vértices de um tetraedro.
\ce{Ca^{2+}}: 8 \ce{F-} nos vértices de um cubo.
\textbf{4.} $a\sqrt3/4 = 546.3 \times 0.4330 = \qty{236.6}{pm}$.
\textbf{5.} $112 + 131 = \qty{243}{pm}$: diferença menor que 3\,\%.
\textbf{6.} Os íons \ce{F-} formam um arranjo cúbico simples de aresta $a/2 =
\qty{273.2}{pm}$, um pouco mais que $2 \times 131 = \qty{262}{pm}$: eles
quase se tocam.
\textbf{7.} $x = 112/131 = 0.85$.
\textbf{8.} $x > 0.732$: coordenação 8.
\textbf{9.} O \omterm{def:b1:crystals-ionic-covalent:structure-type}{tipo cloreto de césio} tem tantos cátions quanto ânions; com
o dobro de ânions, só metade dos cubos de ânions pode conter um cátion.
\textbf{10.} Os 8 \ce{F-} da célula formam um cubo de 8 cubinhos de aresta
$a/2$; os íons \ce{Ca^{2+}} ocupam os centros de 4 deles, alternadamente,
como as casas pretas de um tabuleiro de xadrez tridimensional.
\textbf{11.} $C = \frac43\pi(4 \times 112^3 + 8 \times 131^3)/546.3^3 = 0.61$.
\textbf{12.} \ce{O^{2-}} na \omterm{def:b1:crystals-metals:lattice}{rede} CFC (vértices e centros de face),
\ce{Li+} nos 8 \omterm{def:b1:crystals-metals:site}{sítios tetraédricos}.
\textbf{13.} $a\sqrt3/4 = \qty{200.0}{pm}$; $59 + 142 = \qty{201}{pm}$.
\textbf{14.} $\rho = 4 \times 29.88/(6.022 \times 10^{23} \times (4.619
\times 10^{-8})^3) = \qty{2.01}{g/cm^3}$, o valor medido.
\textbf{15.} A blenda: o diamante é uma blenda com carbono nos dois tipos
de posições.
\textbf{16.} $a\sqrt3/4 = \qty{154.5}{pm}$; $\rho = 8 \times
12.01/(6.022 \times 10^{23} \times (3.567 \times 10^{-8})^3) =
\qty{3.52}{g/cm^3}$.
\textbf{17.} No diamante, só metade dos \omterm{def:b1:crystals-metals:site}{sítios tetraédricos} está ocupada,
e por átomos tão grandes quanto os da \omterm{def:b1:crystals-metals:lattice}{rede}, que só se tocam ao longo das
ligações: $C = 0.34$. Na fluorita, todos os sítios estão ocupados e os
cátions pequenos preenchem os vazios entre ânions grandes: $C = 0.61$.
\textbf{18.} $40.08 + 2 \times 19.00 = \qty{78.08}{g/mol}$.
\textbf{19.} $4 \times 78.08/(6.022 \times 10^{23}) = \qty{5.186e-22}{g}$.
\textbf{20.} $(5.463 \times 10^{-8})^3 = \qty{1.630e-22}{cm^3}$.
\textbf{21.} Um \ce{F-} faltando a cada cem células retira $19.00/(100
\times 312.3)$ da massa, 0.06\,\%: invisível com três algarismos (e um
\ce{Ca^{2+}} também precisa sair, senão o \omterm{def:b1:crystals-metals:crystal}{cristal} não seria neutro).
\textbf{22.} $\rho = 5.186 \times 10^{-22}/1.630 \times 10^{-22} =
\textbf{\qty{3.18}{g/cm^3}}$, exatamente a massa específica medida: o
modelo da célula, construído a partir de uma única medida de $a$ por
raios X, explica a massa do mineral.
\end{solution}
